\chapter{Proposed Methodology and Implementation}
\label{Chapter 3}

This chapter is the long one. I want to walk through how the team actually works at Ignosis, then show the platform from a thousand feet, then go subsystem by subsystem through everything I touched. The chapter is feature-driven on purpose. The platform's brain is a tenant-configurable workflow engine, and basically every other piece -- the campaign module, the AI voice integration, the WhatsApp service, the speed dialer -- is either pushing work into that engine or doing something with the actions it spits out.

\section{Development Methodology}
Engineering at Ignosis is plan-first, ship-small, very review-heavy. I didn't appreciate what that meant in my first week. I wrote what I thought was a small refactor of the tenant-context propagation across async boundaries (my second ticket), pushed it to a feature branch, and got back nine comments, two of which said ``do you want to write a one-pager first and we'll talk through it''. That was the lesson. Every task starts as a short markdown plan attached to the ticket. The engineering lead reads it. You go back and forth in the ticket for a day, maybe two. Only after that does any code get written.

The rest of the workflow is conventional. A feature branch goes up from develop. Conventional commits are mandatory and the CI enforces them. PRs to develop are rebased, not squashed, on purpose, so the per-commit decision trail survives the merge -- this surprised me at first, but it's saved me at least twice when I needed to figure out why some line of code existed. The release branch is on semantic versioning and is cut off main. Tests sit next to the code they exercise. There is a unit-test framework on the backend and a parallel test runner on the frontend. CI fails any build where the static-analysis tool flags a new bug, a new code smell above medium severity, or a new vulnerability of any severity, and there is a fifty-percent coverage floor on changed lines. (I learned about the coverage floor the way you'd expect, which was by failing a PR on a Friday evening.)

\section{High-Level Architecture}
The platform is a layered web app. At the front are three client surfaces: the agent dashboard in the browser, a desktop client that runs the WhatsApp workers, and an Android SMS sender app. Behind those sit the orchestrator backend and a smaller automation service. The data layer is a relational database (cases, campaigns, jobs), a search index for free-text, an analytics event store for reporting, and an in-memory queue with locks. Underneath the data layer the platform talks to three families of external services: the internal voice AI platform plus a telephony gateway for actually placing calls, the WhatsApp protocol through a third-party multi-device library, and the cellular SMS network through whatever the Android phone has built in. Figure~\ref{fig:arch} captures all of that in one picture.

\begin{figure}[H]
\centering
\begin{tikzpicture}[
  font=\footnotesize,
  node distance=0.45cm and 0.4cm,
  box/.style={rectangle, draw, rounded corners=2pt, minimum height=0.85cm, minimum width=2.45cm, align=center, fill=blue!10},
  apibox/.style={box, fill=orange!18},
  databox/.style={box, fill=green!18},
  combox/.style={box, fill=red!15},
  custbox/.style={box, fill=yellow!25, minimum width=4cm},
  lbl/.style={font=\footnotesize\itshape\bfseries, anchor=west},
  arr/.style={-Latex, thick}
]
\node[box] (react) {Agent dashboard\\(browser app)};
\node[box, right=of react] (electron) {Desktop client\\(WhatsApp workers)};
\node[box, right=of electron] (mobile) {Android SMS\\sender app};

\node[apibox, below=1.1cm of react] (kc) {Identity\\provider};
\node[apibox, right=of kc] (spring) {Orchestrator\\backend};
\node[apibox, right=of spring] (hono) {Automation\\API};

\node[databox, below=1.0cm of kc] (pg) {Relational DB\\cases, jobs};
\node[databox, right=of pg] (solr) {Search\\index};
\node[databox, right=of solr] (ch) {Analytics\\event store};
\node[databox, right=of ch] (rds) {In-memory\\queue \& locks};

\node[combox, below=1.0cm of pg] (voice) {Internal voice AI\\+ telephony};
\node[combox, right=of voice] (wa) {WhatsApp\\(via SIM pool)};
\node[combox, right=of wa] (sms) {SMS\\(Android SIM)};

\node[custbox, below=0.7cm of wa] (cust) {Customer (delinquent case)};

\node[lbl] at ($(react.west)+(-2.7,0)$) {Client};
\node[lbl] at ($(kc.west)+(-2.7,0)$) {API};
\node[lbl] at ($(pg.west)+(-2.7,0)$) {Data};
\node[lbl] at ($(voice.west)+(-2.7,0)$) {Channels};

\begin{scope}[on background layer]
  \node[fit=(react)(mobile), draw, dashed, inner sep=4pt] {};
  \node[fit=(kc)(hono), draw, dashed, inner sep=4pt] {};
  \node[fit=(pg)(rds), draw, dashed, inner sep=4pt] {};
  \node[fit=(voice)(sms), draw, dashed, inner sep=4pt] {};
\end{scope}

\draw[arr] (react) -- (spring);
\draw[arr] (electron) -- (hono);
\draw[arr] (mobile.south) to[bend right=25] (hono.north east);
\draw[arr] (spring) -- (pg);
\draw[arr] (spring) -- (solr);
\draw[arr] (spring) -- (ch);
\draw[arr] (spring) -- (rds);
\draw[arr] (hono) -- (rds);
\draw[arr] (spring) -- (voice);
\draw[arr] (spring) -- (wa);
\draw[arr] (hono) -- (wa);
\draw[arr] (hono) -- (sms);
\draw[arr] (voice) -- (cust);
\draw[arr] (wa) -- (cust);
\draw[arr] (sms) -- (cust);
\draw[arr, dashed] (cust.east) to[bend right=40] node[midway, fill=white, font=\tiny] {webhook / SMS reply} (spring.east);
\end{tikzpicture}
\caption{High-level architecture. Solid arrows show forward dispatch; the dashed arrow shows the return path used by voice-agent webhooks and Android SMS-reply ingestion.}
\label{fig:arch}
\end{figure}

It took me longer than I'd like to admit to grasp how this all fits together. I spent most of my first two weeks just reading the orchestrator and chasing where things called what. The choice of a single multi-tenant monolith for the backend instead of a fleet of microservices was made by the team well before I joined, and it was deliberate. The team is small (six engineers, growing slowly), the deployment volume is modest, and the operational cost of running microservices wasn't worth what you'd get back. The boundaries inside the orchestrator are kept clean enough that, if the volume ever justifies a split, the rewrite would not be a rewrite.

\section{Case Management}
\label{sec:case-mgmt}
The unit of work everywhere in the platform is the {\em case}. A case is one delinquent loan account allocated to the tenant for the current cycle. It carries the customer's contact info, the loan account number, the principal and EMI amount, the DPD, the assigned bucket, the agent or AI track it's been allocated to, the workflow node it currently sits on, and a running log of every interaction that has ever happened against it.

The case listing screen, which was my first real frontend ticket, is a server-side-paginated table of the cases assigned to the logged-in agent. Columns are customer name, loan account number, DPD, overdue amount, last contact attempt and current status. Filters get evaluated on the backend; free-text search hits a dedicated index. The case detail screen pulls together the contact info, loan and payment history, every prior call transcript and outcome, WhatsApp and SMS threads, AA-derived signals when those are available, and the case's current position in the workflow. The first version of that page made one big sequential database round-trip and took roughly six seconds to render against a tenant with a busy timeline. Mr. Shah's review of my PR was three words: ``progressive loading please''. So now the core info shows up first, and the interaction-heavy panels stream in behind it in parallel. That single change brought the page to under a second on the same tenant.

\section{The Workflow Engine}
\label{sec:workflow}
The single thing that varies most between tenants is collection policy. A private bank wants a calm tele-calling sequence with WhatsApp reminders sprinkled in between. An NBFC wants an aggressive AI voice campaign in the first ninety-six hours and an escalation to a human only if the AI gives up. Hard-coding any of this would force a code change and a deployment every time a tenant tweaked its policy. The earliest version of the orchestrator did exactly this, and it was painful enough that the team replaced it before I joined.

So the platform stores the policy as data instead. Every tenant has one or more workflows. A workflow is a directed graph: the nodes are concrete actions (initiate AI call, send WhatsApp reminder, send payment link, schedule agent callback, wait for payment, escalate to human), and the transitions are simple boolean conditions on the case's current variables. A transition reads like ``advance to the wait-for-payment node if the customer promised to pay in the next three days; otherwise fall through''. The engine fetches the current node, runs its action (synchronously if it's cheap, asynchronously via a scheduled job if it's not), drops the outcome into the case's variable context, walks the outgoing transitions in declaration order, and parks the case on the next node. If none of the transitions match, the case goes into an error state and QA picks it up. That last bit -- the QA error state -- exists because we had a handful of badly-authored workflows that silently looped in the first month of one of the tenants. Visible failure beats invisible failure here.

Every workflow action that touches an external system runs asynchronously. The engine writes a job into a database-backed queue, and a scheduler polls for due jobs, takes a row-level lock with a short time-to-live and runs the executor. Why a database-backed queue and not an in-memory one? Visibility. When a campaign is going sideways at 3 AM the on-call engineer can query the job state directly and see exactly where each case is stuck. The trade-off is throughput. At the campaign volumes the platform actually sees in production, throughput hasn't been a problem.

\section{The Campaign Module}
\label{sec:campaign}
A campaign is the platform's bulk-work unit: ``run this workflow against this batch of cases, inside this calling window''. The creation form takes a campaign name, a priority, a target collection amount, an assignment rule (round-robin, weighted, manual) and a calling window. The CSV upload accepts files up to fifty megabytes. The CSV-processing service parses the file, validates each row against a small set of rules (non-blank phone in international format, non-blank name, positive due amount, valid loan account number, recognised bucket, in-range DPD), and returns a preview with per-row diagnostics. The campaign manager looks over the preview, fixes anything obvious, and submits the campaign for execution.

Two things about this flow that are not obvious. First, the validator rejects rows but does not reject the file: the manager almost always has a handful of bad rows in a 4,000-row CSV (a phone number missing the country code, a bucket that doesn't match the DPD), and stopping the entire upload on those would make the tool unusable. So bad rows are flagged in the preview and skipped on submission, and a downloadable report shows what got dropped. Second, one of the tenants kept uploading CSVs with a BOM marker at the start of the first header, which broke the parser silently for the first two weeks of that onboarding. That bug is fixed (the parser now handles the BOM), but it's why the preview now also includes a ``did we recognise your columns'' line.

On submission, the backend persists the campaign and one case record per valid row, and the campaign moves into a scheduled or ongoing state. Once it's ongoing, the scheduler picks it up on the next tick and starts enqueueing workflow-advancement jobs for its rows. From here on, the workflow engine drives every case independently. The campaign object itself keeps track of progress (cases by state, total amount collected, calls attempted versus connected) and pushes those numbers to the dashboard over a streaming channel. A campaign can be paused, resumed or marked complete. The biggest campaign we have run through the platform so far was 41,200 cases. It took the scheduler about seven minutes to get all the initial advancement jobs queued. That run, and the question of whether the queue could keep up, was the first time I had to actually think about throughput on the platform.

\section{AI Voice Calling}
\label{sec:voice}
This is the headline feature of the platform. From an architectural standpoint, an AI call is just one of the action types the workflow engine knows how to run. When the engine reaches an AI call node for a case, it schedules a job. The executor reads the tenant's configured voice-agent identifier, pulls the next phone number from the tenant's pool of caller-line identifiers, builds the call request from the case context (customer name, due amount, DPD, AA signals if available), and hands the request off to the internal voice AI platform. From here on the voice platform owns the call until it terminates.

The voice platform places the call through the telephony gateway, runs the conversation under the tenant's prompt template, and fires a callback to the orchestrator when the call ends. The callback carries the disposition, a structured extraction of variables (promise-to-pay date, dispute reason, callback request, language detected), a recording URL and the per-turn latency. The orchestrator parses it, persists the outcome onto the case, cancels any defensive status-poll job that was scheduled as a safety net, emits metrics for the call, and re-enters the workflow engine for that case so the workflow advances on the same callback request.

The first time I listened to one of these recordings end to end, I was sitting in the team area on a Tuesday and a customer in Lucknow had agreed to pay on the 18th. The AI never raised its voice once. I think that's when I actually understood why the team was so invested in this subsystem.

Figure~\ref{fig:voice-seq} lays the full end-to-end interaction out as a sequence diagram, in the order in which the orchestrator, voice platform, telephony gateway, customer and workflow engine exchange messages. The dashed return arrows are the inbound callback path that re-enters the engine on the same request.

\begin{figure}[H]
\centering
\begin{tikzpicture}[
  font=\footnotesize,
  lifeline/.style={dashed, gray!70},
  actor/.style={rectangle, draw, fill=blue!12, minimum width=1.6cm, minimum height=0.55cm, align=center, font=\footnotesize\bfseries},
  msg/.style={-Latex, thick},
  ret/.style={-Latex, thick, dashed},
  note/.style={font=\tiny\itshape, fill=yellow!18, draw=gray!40, inner sep=2pt, rounded corners=1pt}
]
\def\xa{0}
\def\xb{2.7}
\def\xc{5.4}
\def\xd{8.1}
\def\xe{10.8}
\def\ytop{0}
\def\ybot{-10.8}

\node[actor] (orch) at (\xa,\ytop) {Orchestrator};
\node[actor] (eng)  at (\xb,\ytop) {Workflow\\engine};
\node[actor] (bolna) at (\xc,\ytop) {Voice AI\\platform};
\node[actor] (exo)  at (\xd,\ytop) {Telephony\\gateway};
\node[actor] (cust) at (\xe,\ytop) {Customer};

\foreach \x in {\xa,\xb,\xc,\xd,\xe} {
  \draw[lifeline] (\x,-0.3) -- (\x,\ybot);
}

\draw[msg] (\xa,-0.8)  -- node[above, font=\tiny] {advance(case)} (\xb,-0.8);
\draw[msg] (\xb,-1.4)  -- node[above, font=\tiny] {schedule call-initiate job} (\xa,-1.4);
\node[note, anchor=west] at (\xa+0.05,-1.95) {also schedules 5-min poll job (defensive)};
\draw[msg] (\xa,-2.6)  -- node[above, font=\tiny] {initiate call (case context, from-number)} (\xc,-2.6);
\draw[msg] (\xc,-3.2)  -- node[above, font=\tiny] {dial} (\xd,-3.2);
\draw[msg] (\xd,-3.8)  -- node[above, font=\tiny] {ring} (\xe,-3.8);
\draw[ret] (\xe,-4.4)  -- node[above, font=\tiny] {answer} (\xd,-4.4);
\draw[ret] (\xd,-5.0)  -- node[above, font=\tiny] {audio stream} (\xc,-5.0);
\node[note, anchor=west] at (\xc+0.05,-5.55) {ASR + LLM + TTS turn loop ($\sim$700~ms median)};
\draw[msg] (\xc,-6.2)  -- node[above, font=\tiny] {audio stream} (\xd,-6.2);
\draw[msg] (\xd,-6.8)  -- node[above, font=\tiny] {speech} (\xe,-6.8);
\draw[ret] (\xe,-7.4)  -- node[above, font=\tiny] {customer hangs up} (\xd,-7.4);
\draw[ret] (\xd,-8.0)  -- node[above, font=\tiny] {call ended} (\xc,-8.0);
\draw[ret] (\xc,-8.6)  -- node[above, font=\tiny] {callback: disposition, transcript, recording URL} (\xa,-8.6);
\draw[msg] (\xa,-9.2)  -- node[above, font=\tiny] {persist outcome, cancel poll job} (\xb,-9.2);
\draw[msg] (\xb,-9.8)  -- node[above, font=\tiny] {evaluate transitions, advance node} (\xa,-9.8);
\node[note, anchor=west] at (\xa+0.05,-10.4) {workflow advances on the same callback request};
\end{tikzpicture}
\caption{Sequence diagram for an AI voice call. Solid arrows are forward dispatch; dashed arrows are returns (audio streams and the terminal callback). The defensive 5-minute poll job (not drawn) is cancelled when the callback arrives first.}
\label{fig:voice-seq}
\end{figure}

The speed dialer is the user-facing twin of the AI calling path. While an agent is working on a campaign, the dashboard shows the active case in a panel. The agent talks to the customer, fills out a short disposition form and hits submit. The instant the disposition lands at the backend, the workflow engine advances the case, the case enters a terminal-for-this-agent state, and the next pending case is pushed to the dashboard over the streaming channel. From the agent's seat, the next customer is already on screen as the previous form is closing. The agency floor manager I demoed this to in March told me her agents had stopped using sticky notes for the next number to dial. I think that's the cleanest user feedback any of my work has gotten this year.

\section{WhatsApp Automation}
\label{sec:whatsapp}
The WhatsApp automation subsystem lives in a separate automation toolkit, deliberately kept apart from the orchestrator for operational reasons. WhatsApp sessions are long-lived, stateful connections. They need physical phones for the initial QR pairing. They have to survive infrastructure restarts. None of those things sits comfortably in a stateless backend service. The toolkit is built on a third-party multi-device protocol library that drives a pool of consumer SIMs in software. Choosing this library over the official business API was an easy call. The business API has template approvals that drag on for days, per-message pricing, and customers reliably respond better to messages from real consumer numbers than from obviously-branded business accounts.

The first version of this subsystem ran on a single phone. I had it as a CLI tool that took a JSON file of contacts and sent messages out one at a time. It was fine for ten messages. It fell over at a hundred. The rewrite came two weeks later as a service driving a pool of phones with a shared queue. The third version, which is what's in production now, added the preferred-sender routing, the per-sender daily caps, the warm-up schedule and a real block-recovery path. I'd love to say the three versions were planned. They were not. The first was prototyped without enough thought. The second was rewritten on a Sunday because I'd promised myself the weekend off and then I couldn't stop thinking about a bug. The third was the one that actually had a design document, written after Mr. Shah caught me planning the second rewrite on a whiteboard and said, gently, that this was the kind of thing we wrote down.

A campaign that uses WhatsApp is bound to a pool of five to fifteen sender phones. When the campaign starts, the scheduler computes a proportional allocation of contacts across the pool, capping each sender at its own daily limit. A shared contact queue carries a {\em preferred sender} field so that a contact can be soft-routed to a specific phone, usually one that already has an existing thread with the customer. If the preferred sender is offline or banned, the contact falls back to whichever worker picks it up next. That makes the system resilient when individual senders die mid-cycle, which they routinely do. We lost three senders in one afternoon during a campaign in mid-March and the system reassigned the orphans cleanly. I was holding my breath the whole time.

A real chunk of the engineering work on this subsystem goes into staying ahead of spam-detection heuristics. Each sender has a daily cap. The inter-message delay is jittered by forty percent around a per-sender base interval. Batches of three to twelve messages are followed by short breaks; every fifth batch triggers a longer break. Workers honour an active-hours window. Every outgoing message body is varied through inline alternates and a zero-width-space inserted at a random character position. The message-variation idea came out of a Stack Overflow thread I'd half-read on a Tuesday morning. Block detection is reactive but fast: when the protocol library signals a ban, the sender is pulled out of rotation, the suspected ban is logged, and the orphaned contacts in the shared queue are reassigned to the remaining workers. Transient connection-class errors get handled separately with exponential backoff. A workflow node that sends a WhatsApp reminder posts a job to the automation service, the service runs the send, waits for the protocol-level acknowledgement, and returns the delivery outcome to the orchestrator. The workflow then advances on the back of that outcome.

\section{SMS via Android}
The SMS fallback channel is a small Android app that runs on a phone with an active consumer SIM. It calls into the device's built-in SMS interface to send and read messages. Outbound sends respect a per-message delay, a batch size and a per-batch pause, all to keep the SIM operator from flagging the device. The numbers we use here -- twelve seconds between messages, batches of twenty, a forty-five-second gap between batches -- came out of two weeks of careful trial and error on a single Redmi we keep in the office. An evening background task reads new SMS messages off the on-device provider, bundles them up and uploads them to the orchestrator so the agent sees every reply alongside the rest of the case timeline. From the workflow engine's point of view, an SMS reminder node is exactly the same shape as a WhatsApp reminder node. It produces an outcome. The outcome drives the next transition.

\section{The Desktop Application}
The desktop application is the operational vehicle that hands the WhatsApp workers to the engineering team and to a small set of agency-side operators. On launch it boots the embedded automation server, serves the dashboard locally and opens a window into it. It has two main jobs. One is running WhatsApp campaigns from the operator's own machine with locally-paired SIMs. The other is scraping the chat history of the paired phones at the end of each day so the case dashboard can show every WhatsApp conversation alongside the call history.

The first build hard-coded the local server's port. This worked on three of the four machines we tested it on. On the fourth, it clashed with another process and the app showed a one-line cryptic error that didn't actually say anything about a port conflict. I spent the better part of an afternoon on a video call with the operator before working out what had happened. The port-finding logic that's now in the second build should have been in the first.

\section{Data Analysis on Customer Records}
Very early on -- in my fourth week, I think -- I ran a one-off analysis on a sample of roughly forty thousand customer records exported from one of the lender tenants. The dataset arrived as a 220 MB CSV with about twenty-six columns and a number of obvious quality issues that I had to clean up before anything sensible would come out. I sat with it for two days. Three patterns came out with practical implications.

First, customers who had paid on the first attempt in any of the previous three cycles had a ninety-percent probability of paying on the first attempt in the next cycle. These are best served by a single SMS reminder rather than a full call campaign, which saves both the agency time and the lender money. Second, customers whose detected salary credit landed in the first three days of the month had a noticeably higher response rate to outreach attempted on the fourth or fifth. The size of that effect was not huge but it was very consistent across tenants. Third, there is a long tail of customers with ticket sizes under one thousand rupees and recovery probabilities below ten percent, for which a single SMS is the only touch that makes any economic sense. All three of these findings went into the default workflow templates the platform ships with for new tenants.

\section{Summary of Contributions}
The pieces of the platform I contributed: the case listing and case detail screens on the agent dashboard; the bulk campaign module, including the creation form and the CSV upload pipeline; the speed dialer and the integration of AI voice calling into the workflow engine; the WhatsApp automation, including session management, the sender pool scheduler and the anti-ban heuristics; multi-device rotation, block detection and message variation; the desktop application; the Android SMS sender app with its background sync; the voice-platform callback integration on the workflow path; and the one-off analysis of the customer records. Each of these is either producing work the workflow engine consumes, consuming work the workflow engine produces, or shortening how long a human agent spends on cases the workflow has already prioritised. Every number in Chapter~\ref{Chapter 5} traces back to that workflow-centric structure.
